\documentclass[conference]{IEEEtran}
\IEEEoverridecommandlockouts

\usepackage[utf8]{inputenc}
\usepackage[T1]{fontenc}
\usepackage[turkish]{babel}
\usepackage{cite}
\usepackage{amsmath,amssymb,amsfonts}
\usepackage{algorithmic}
\usepackage{graphicx}
\usepackage{textcomp}
\usepackage{xcolor}
\usepackage{booktabs}
\usepackage{listings}
\usepackage{hyperref}
\usepackage{float}

\lstset{
  basicstyle=\ttfamily\footnotesize,
  breaklines=true,
  frame=single,
  numbers=left,
  numberstyle=\tiny,
  keywordstyle=\color{blue},
  commentstyle=\color{gray},
  stringstyle=\color{red},
  language=Python
}

\begin{document}

% ── KAPAK SAYFASI ──────────────────────────────────────────────────────────────
\begin{titlepage}
  \centering
  \vspace*{2cm}

  {\Large \textbf{OSTİM TEKNİK ÜNİVERSİTESİ}\\[0.3cm]
  Yapay Zeka Mühendisliği Bölümü}

  \vspace{1.5cm}
  \rule{\linewidth}{0.5mm}\\[0.4cm]

  {\LARGE \textbf{xPass++: Graf Sinir Ağı ile Futbolda\\[0.2cm]
  Pas ve Şut Karar Kalitesi Tahmini}}\\[0.2cm]

  \rule{\linewidth}{0.5mm}\\[1cm]

  {\large \textbf{Derin Öğrenme Dersi Final Ödevi}}

  \vspace{1.5cm}

  \begin{tabular}{ll}
    \textbf{Öğrenci Adı Soyadı:} & Mustafa Necati ÖLMEZ \\[0.3cm]
    \textbf{Öğrenci Numarası:}   & 230212030 \\[0.3cm]
    \textbf{Bölüm:}              & Yapay Zeka Mühendisliği \\[0.3cm]
    \textbf{Danışman:}           & Dr. Murat ŞİMŞEK \\[0.3cm]
    \textbf{Tarih:}              & Mayıs 2026 \\
  \end{tabular}

  \vspace{1.5cm}

  \textbf{GitHub:} \url{https://github.com/mustafaaolmez/xpass-gnn}\\[0.3cm]
  \textbf{Web Arayüzü:} \url{https://mustafaaolmez.github.io/xpass-gnn/}

  \vfill
  {\small Ankara, Türkiye}
\end{titlepage}
% ── KAPAK SONU ─────────────────────────────────────────────────────────────────

\title{xPass++: Graf Sinir Ağı ile Futbolda Pas ve Şut Karar Kalitesi Tahmini}

\author{
  \IEEEauthorblockN{Mustafa Necati ÖLMEZ}
  \IEEEauthorblockA{
    \textit{Yapay Zeka Mühendisliği Bölümü}\\
    \textit{OSTİM Teknik Üniversitesi}\\
    Ankara, Türkiye\\
    230212030\\
    Derin Öğrenme Dersi Final Projesi\\
    Danışman: Dr. Murat ŞİMŞEK
  }
}

\maketitle

% ── ÖZET ────────────────────────────────────────────────────────────────────────
\begin{abstract}
Bu çalışmada, futbolda pas ve şut kararlarının kalitesini ölçen \textbf{xPass++} adlı
bir Graf Sinir Ağı (Graph Neural Network, GNN) sistemi tasarlanmış ve geliştirilmiştir.
FIFA 2022 Dünya Kupası verisi (StatsBomb Open Data) üzerinde, 360$^\circ$ dondurulmuş
kare (freeze-frame) girdileri ve xT-delta denetimli hedef kullanılarak eğitilen model,
sahadaki oyuncu etkileşimlerini graf yapısıyla modellemektedir.
Nihai \textbf{Karar Kalitesi Skoru (Decision Quality Score, DQS)} = kapalı-form xT-temeli
+ Graf Dikkat Ağı (Graph Attention Network, GAT) tarafından tahmin edilen artık değer
formülüyle hesaplanmaktadır.
Test bölümünde ($n = 7.394$) xT-temeli baseline'a kıyasla MAE'de \textbf{\%10.5},
MSE'de \textbf{\%32} iyileşme elde edilmiştir. Model 3 GAT katmanı, 4 dikkat kafası ve
128 gizli boyut içermekte; toplam $\approx$211.000 parametreye sahip bulunmaktadır.
\end{abstract}

\begin{IEEEkeywords}
Graf Sinir Ağı, Graph Attention Network, Futbol Analitigi, Karar Kalitesi, xPass, xT,
DQS, StatsBomb, PyTorch Geometric
\end{IEEEkeywords}

% ── 1. GİRİŞ ────────────────────────────────────────────────────────────────────
\section{Giriş}

Futbol analitiği, son on yılda istatistiksel modellemenin en hızlı geliştiği alanlardan
biri hâline gelmiştir. Beklenen gol (xG), beklenen pas (xPass) ve tehdit transferi (xT)
gibi metrikler kulüpler tarafından yaygın biçimde kullanılmaktadır~\cite{fernandez2019decomposing}.
Ne var ki mevcut metrikler çoğunlukla tek bir eylemi izole ederek değerlendirir;
oyuncunun o an sahada oluşturduğu bağlamı bütünsel olarak ele almaz.

Bir oyuncu pas verirken aslında bir karar alır: o konumdan, o zamanda, etraftaki
oyuncuların dizilimiyle mümkün olan en iyi seçeneği seçmek. Klasik istatistikler
``pas başarılı mı?'' diye sorar. \textbf{xPass++} ise ``bu pas ne kadar kaliteli bir
karardı?'' sorusunu yanıtlamaya çalışır.

Bu soruyu yanıtlamak için Graf Sinir Ağları (GNN) doğal bir çerçeve sunar: sahadaki
oyuncular düğüm, aralarındaki mekânsal ilişkiler ise kenar olarak modellenir. Graf
Dikkat Ağı (GAT)~\cite{velickovic2018gat}, her düğümün komşularına ne kadar dikkat
etmesi gerektiğini öğrenerek futbol oyununun ilişkisel yapısını kapsar.

Bu çalışmanın temel katkıları şunlardır:
\begin{enumerate}
  \item Futbol eylemlerini yönlü-ağırlıklı graf olarak modelleyen bir veri hattı.
  \item Kapalı-form xT skoru ile GAT artık tahminini birleştiren yorumlanabilir DQS
  çerçevesi.
  \item FIFA 2022 Dünya Kupası verisinde (64 maç, 73.946 eylem) baseline'a kıyasla
  MAE'de \%10.5 iyileşme.
\end{enumerate}

% ── 2. İLGİLİ ÇALIŞMALAR ────────────────────────────────────────────────────────
\section{İlgili Çalışmalar}

\subsection{Futbol Analitiğinde İstatistiksel Modeller}

Beklenen gol (xG) modelleri~\cite{rathke2017expected} şut kalitesini şutun konumu,
vücut parçası ve oyun fazı gibi özniteliklerle tahmin eder. Fernandez ve ark.~\cite{fernandez2019decomposing}
tehdit transferi (xT) kavramını formüle ederek saha bölgelerinin değerini belirli bir
Markov süreci çerçevesinde hesaplamıştır. Bu kapalı-form yaklaşımlar yorumlanabilir
olmakla birlikte oyuncu etkileşimlerini göz ardı eder.

\subsection{Graf Sinir Ağları ile Spor Analizi}

Alciaturi ve ark.~\cite{alciaturi2023graph} futbol baskı (pressing) eylemlerini GNN
ile modellemiş; oyuncu konumlarını düğüm olarak kullanmıştır. Power ve ark.~\cite{power2017optimising}
pas seçimini optimizasyon problemi olarak çerçevelemiştir. Bununla birlikte, karar
kalitesini doğrudan bir artık (residual) modelleme yaklaşımıyla tahmin eden bir GNN
çalışması literatürde bulunmamaktadır.

\subsection{Graf Dikkat Ağları}

Velicković ve ark.~\cite{velickovic2018gat} tarafından önerilen GAT, her komşu
düğüme öğrenilmiş dikkat ağırlıkları atayarak standart grafik evrişimini genişletir.
Çok başlıklı dikkat mekanizması farklı ilişki türlerini eş zamanlı olarak
öğrenebilmekte; bu özellik oyuncu etkileşimlerinin çok boyutlu yapısını modellemeye
uygun bir çerçeve sunmaktadır~\cite{hamilton2020graph}.

% ── 3. VERİ KÜMESİ ──────────────────────────────────────────────────────────────
\section{Veri Kümesi}

\subsection{StatsBomb Open Data}

Veri kaynağı olarak StatsBomb'un ücretsiz açık veri deposu~\cite{statsbomb2023open}
kullanılmıştır. FIFA 2022 Dünya Kupası maçlarını kapsayan bu veri seti, her eylem için
360$^\circ$ dondurulmuş kare (freeze-frame) içermektedir. Bu kareler, eylemin
gerçekleştiği anda sahadaki tüm oyuncuların konumunu ve takım bilgisini içerir.

\begin{table}[H]
  \centering
  \caption{Veri Kümesi İstatistikleri}
  \label{tab:dataset}
  \begin{tabular}{lr}
    \toprule
    \textbf{Özellik} & \textbf{Değer} \\
    \midrule
    Turnuva & FIFA 2022 Dünya Kupası \\
    Maç Sayısı & 64 \\
    Toplam Graf Sayısı & 73.946 \\
    Eylem Türleri & Pas + Şut \\
    Düğüm Öznitelik Boyutu & 16 \\
    Kenar Öznitelik Boyutu & 8 \\
    \midrule
    Eğitim Bölümü & \%70 (51.762) \\
    Doğrulama Bölümü & \%15 (11.092) \\
    Test Bölümü & \%15 (11.092) \\
    \bottomrule
  \end{tabular}
\end{table}

\subsection{Graf Yapısı}

Her eylem için bir graf $G = (V, E)$ oluşturulur. Düğüm kümesi $V$, eylemin
gerçekleştiği andaki tüm görünür oyuncuları (ortalama 17-22 oyuncu) temsil eder.
Kenar kümesi $E$, belirli bir mesafe eşiği ($d_{thr} = 10$ m) içindeki oyuncu
çiftlerini birbirine bağlar. Eylem sahibi düğüm ``merkez düğüm'' olarak
etiketlenir ve DQS tahmini bu düğüm için yapılır.

Düğüm öznitelikleri şunlardır: normalize edilmiş $x$/$y$ konumu, hız vektörü,
takım kimliği (ikili), topu elinde bulundurma durumu, kaleciye uzaklık ve kaleye
uzaklık. Kenar öznitelikleri: oyuncular arası mesafe, görüş açısı ve takım aynılığı
durumunu içermektedir.

\subsection{Hedef Etiketler (xT-delta)}

Hedef etiket, bir eylemin xT değerini nasıl değiştirdiğini ölçen $\Delta xT$
(xT-delta) değeridir:

\begin{equation}
  \Delta xT_i = xT(\text{konum}_{\text{sonra}}) - xT(\text{konum}_{\text{önce}})
\end{equation}

Bu değer \texttt{produce\_labels.py} betiği aracılığıyla hesaplanarak
\texttt{wc2022\_xt\_labels.parquet} dosyasına kaydedilmiştir.

% ── 4. YÖNTEM ───────────────────────────────────────────────────────────────────
\section{Yöntem}

\subsection{DQS Çerçevesi}

xPass++ iki bileşenin toplamı olarak tanımlanır:

\begin{equation}
  \text{DQS}_i = xT_{\text{kapalı-form}}(s_i) + \hat{r}_i
  \label{eq:dqs}
\end{equation}

Burada $xT_{\text{kapalı-form}}(s_i)$ konum tabanlı kapalı-form xT tahmini,
$\hat{r}_i$ ise GAT'ın tahmin ettiği artık (residual) değerdir. Bu ayrıştırma
sayesinde model karar açıklanamaz bir kara kutu olmaktan çıkmakta; temel xT skoru
yorumlanabilirliği korurken GAT bağlamsal iyileştirmeyi sağlamaktadır.

\subsection{Model Mimarisi}

xPass++ GAT modeli üç dikkat katmanından oluşmaktadır. Her katmanda çok başlıklı
dikkat mekanizması uygulanmakta; katmanlar arası ELU aktivasyonu ve Dropout
kullanılmaktadır.

\begin{equation}
  \mathbf{h}_i^{(l+1)} = \text{ELU}\!\left(\sum_{j \in \mathcal{N}(i)} \alpha_{ij}^{(l)} \mathbf{W}^{(l)} \mathbf{h}_j^{(l)}\right)
\end{equation}

Dikkat katsayıları $\alpha_{ij}$ şu şekilde hesaplanır:

\begin{equation}
  \alpha_{ij} = \frac{\exp\!\bigl(\text{LeakyReLU}(\mathbf{a}^T [\mathbf{W}\mathbf{h}_i \| \mathbf{W}\mathbf{h}_j])\bigr)}{\sum_{k \in \mathcal{N}(i)} \exp\!\bigl(\text{LeakyReLU}(\mathbf{a}^T [\mathbf{W}\mathbf{h}_i \| \mathbf{W}\mathbf{h}_k])\bigr)}
\end{equation}

\begin{table}[H]
  \centering
  \caption{Model Mimarisi Özeti}
  \label{tab:arch}
  \begin{tabular}{ll}
    \toprule
    \textbf{Bileşen} & \textbf{Değer} \\
    \midrule
    Mimari & Graph Attention Network (GAT) \\
    Katman Sayısı & 3 \\
    Dikkat Kafası & 4 (çok başlıklı) \\
    Gizli Boyut & 128 \\
    Aktivasyon & ELU \\
    Dropout & 0.3 \\
    Toplam Parametre & $\approx$211.000 \\
    \bottomrule
  \end{tabular}
\end{table}

\subsection{Model Kodu}

\begin{lstlisting}[caption={xPass++ GAT Mimarisi (PyTorch Geometric)}]
class xPassPlusPlus(torch.nn.Module):
    def __init__(self, in_ch=16,
                 hidden=128, heads=4):
        super().__init__()
        self.gat1 = GATConv(
            in_ch, hidden,
            heads=heads, dropout=0.3)
        self.gat2 = GATConv(
            hidden * heads, hidden,
            heads=heads, dropout=0.3)
        self.gat3 = GATConv(
            hidden * heads, hidden,
            heads=1, dropout=0.3)
        self.residual = Linear(hidden, 1)

    def forward(self, x, edge_index):
        x = F.elu(self.gat1(x, edge_index))
        x = F.elu(self.gat2(x, edge_index))
        x = F.elu(self.gat3(x, edge_index))
        # Artik tahmin: DQS reziduel
        return self.residual(x)
\end{lstlisting}

% ── 5. EĞİTİM SÜRECİ ────────────────────────────────────────────────────────────
\section{Eğitim Süreci}

\subsection{Eğitim Ortamı ve Donanım}

\begin{table}[H]
  \centering
  \caption{Eğitim Donanımı}
  \label{tab:hardware}
  \begin{tabular}{ll}
    \toprule
    \textbf{Bileşen} & \textbf{Özellik} \\
    \midrule
    Geliştirme GPU & NVIDIA RTX 4050 (6 GB VRAM) \\
    Eğitim Ortamı & Yerel + Google Colab \\
    Python & 3.10 \\
    PyTorch & 2.x \\
    PyTorch Geometric & 2.5+ \\
    \bottomrule
  \end{tabular}
\end{table}

\subsection{Hiperparametreler}

\begin{table}[H]
  \centering
  \caption{Nihai Eğitim Hiperparametreleri}
  \label{tab:hyperparams}
  \begin{tabular}{ll}
    \toprule
    \textbf{Hiperparametre} & \textbf{Değer} \\
    \midrule
    Optimizer & AdamW \\
    Öğrenme Hızı & $1 \times 10^{-3}$ \\
    Ağırlık Azalması & $1 \times 10^{-4}$ \\
    Batch Boyutu & 64 graf/adım \\
    Zamanlayıcı & Cosine Annealing \\
    Maksimum Epoch & 30 \\
    Erken Durdurma & Sabır = 15 \\
    Dropout & 0.3 \\
    Kayıp Fonksiyonu & Huber Loss ($\delta = 1.0$) \\
    \bottomrule
  \end{tabular}
\end{table}

\subsection{İki Aşamalı Eğitim}

Model eğitimi iki sprint'te gerçekleştirilmiştir. \textbf{Sprint-2}'de ham xT
etiketleri doğrudan tahmin edilmiştir. \textbf{Sprint-4 (xPass++)} ise artık
modelleme yaklaşımına geçilmiş; GAT yalnızca kapalı-form xT'nin açıklayamadığı
bağlamsal kalıntıyı tahmin edecek biçimde yeniden tasarlanmıştır. Bu tasarım
değişikliği hem MAE hem de MSE üzerinde belirgin iyileşme sağlamıştır.

% ── 6. SONUÇLAR ─────────────────────────────────────────────────────────────────
\section{Sonuçlar}

\subsection{Nicel Değerlendirme}

\begin{table}[H]
  \centering
  \caption{Model Karşılaştırması (Test Kümesi, $n = 7.394$)}
  \label{tab:results}
  \begin{tabular}{lccc}
    \toprule
    \textbf{Model} & \textbf{MAE} & \textbf{MSE} & \textbf{Spearman $\rho$} \\
    \midrule
    xT-Temeli (Baseline) & 0.00790 & 0.000172 & 0.746 \\
    Sprint-2 GAT         & 0.00750 & 0.000145 & 0.738 \\
    \textbf{xPass++ DQS} & \textbf{0.00707} & \textbf{0.000116} & \textbf{0.731} \\
    \midrule
    $\Delta$ (Baseline $\to$ xPass++) & $-10.5\%$ & $-32.6\%$ & $-0.015$ \\
    \bottomrule
  \end{tabular}
\end{table}

GAT artık kafası kapalı-form baseline'a kıyasla MAE'de \%10.5, MSE'de \%32.6
iyileşme sağlamıştır. Spearman $\rho$ değerinin hafif düşmesi, modelin sıralama
tutarlılığından ziyade mutlak hata minimize etmeye odaklandığını göstermektedir;
bu bir denge tercihidir.

\subsection{Saha Isı Haritası}

\texttt{notebooks/demo\_dqs.py} betiği tek bir maç için her saha bölgesinin ortalama
DQS değerini hesaplayarak görsel bir ısı haritası üretmektedir. Beklentiyle örtüşür
biçimde ceza sahası önü ve merkez orta saha bölgeleri yüksek DQS değerleri
sergilemiştir. Bu bulgular \texttt{docs/figures/demo\_dqs\_pitch.png} dosyasında
görselleştirilmiştir.

% ── 7. PROJE SÜRECİ VE KARŞILAŞILAN ZORLUKLAR ─────────────────────────────────
\section{Proje Süreci ve Karşılaşılan Zorluklar}

\textbf{Graf İnşası:} StatsBomb 360 verisi bazı maçlar için eksik freeze-frame içermektedir.
Bu durum çözülmüş; eksik konumlar atlanarak yalnızca tam freeze-frame içeren eylemler
graflanmıştır.

\textbf{Etiket Gürültüsü:} xT-delta değerleri yüksek varyanslı olabilmektedir. Bu
soruna karşı Huber Loss kullanılmış; aykırı değerlerin gradyan güncellemelerine
aşırı katkısı sınırlandırılmıştır.

\textbf{Artık Tasarım:} Ham etiket tahmini (Sprint-2) yetersiz kaldığında, kapalı-form
xT + GAT artık ayrıştırması (Sprint-4) hem model başarımını artırmış hem de
yorumlanabilirliği korumuştur. Bu mimari karar projenin en kritik dönüm noktası
olmuştur.

\textbf{VRAM Kısıtı:} Büyük batch boyutlarında GPU belleği yetersiz kalmıştır.
Gradient accumulation (4 adım) ile etkin batch boyutu artırılmıştır.

% ── 8. GELECEK ÇALIŞMALAR ───────────────────────────────────────────────────────
\section{Gelecek Çalışmalar}

\textbf{Dinamik Graf:} Mevcut model statik freeze-frame anlık görüntüsünü kullanır.
Zaman içindeki oyuncu hareketini modelleyen dinamik GNN mimarisi (örn. STGNN)
karar bağlamını zenginleştirebilir.

\textbf{Heterojen Graf:} Oyuncu ve top düğümlerini farklı türde düğüm olarak
tanımlayan heterojen GNN, takım içi ve takımlar arası etkileşimleri ayrı ayrı
modelleyebilir~\cite{hamilton2020graph}.

\textbf{Daha Büyük Veri:} Tüm StatsBomb maç arşivi ($>$3.000 maç) ile eğitim,
modelin genelleme kapasitesini önemli ölçüde artırabilir.

\textbf{Açıklanabilirlik:} GNNExplainer veya Grad-CAM'in GNN uyarlaması, hangi
oyuncuların kararı en çok etkilediğini görsel olarak açıklayabilir; bu özellik
kulüp analistleri için pratik değer taşır.

\textbf{Gerçek Zamanlı Sistem:} Model hafifliği ($\approx$211k parametre) göz önüne
alındığında, canlı maç akışıyla entegre gerçek zamanlı DQS hesaplama teknik
açıdan uygulanabilir görünmektedir.

% ── 9. SONUÇ VE DEĞERLENDİRME ──────────────────────────────────────────────────
\section{Sonuç ve Değerlendirme}

Bu proje bana Graf Sinir Ağlarının gerçek dünya verilerine uygulanması konusunda
derin bir pratik deneyim kazandırdı. Futbol verisi, yapısal olarak oldukça
gürültülüdür; eksik değerler, tutarsız etiketler ve dinamik sahne yapısı bütünlüklü
bir veri hattı tasarımı gerektirdi.

Ham etiket tahmini (Sprint-2) başarısız olmasa da sınırlı kalırken, kapalı-form
xT + GAT artık ayrıştırması (xPass++) hem başarımı artırdı hem de yorumlanabilirliği
korudu. Bu ``modeli karmaşıklaştırmak yerine doğru problemi çöz'' prensibi benim
için en değerli öğrenme çıktısı oldu.

PyTorch Geometric ile gerçek bir GNN hattı kurmanın zorlukları --- graf mini-batch
oluşturma, heterojen kenar türleri, attention görselleştirme --- teorik bilgiyi
uygulama becerisiyle pekiştirdi. Proje, hem akademik düzeyde makul sonuçlar üretti
hem de GitHub Pages aracılığıyla herkesin erişebileceği bir web arayüzüyle
paylaşılabilir hâle getirildi.

% ── KAYNAKÇA ────────────────────────────────────────────────────────────────────
\begin{thebibliography}{00}

\bibitem{fernandez2019decomposing}
J. Fernandez, L. Bornn ve D. Cervone,
``Decomposing the Immeasurable Sport: A deep learning expected possession value
framework for soccer,''
\textit{MIT Sloan Sports Analytics Conference}, 2019.

\bibitem{velickovic2018gat}
P. Velicković, G. Cucurull, A. Casanova, A. Romero, P. Liò ve Y. Bengio,
``Graph Attention Networks,''
\textit{arXiv preprint arXiv:1710.10903}, 2018.

\bibitem{alciaturi2023graph}
P. Alciaturi, J. Robson ve C. Dick,
``Graph Neural Networks for football pressing analysis,''
\textit{arXiv preprint arXiv:2306.10055}, 2023.

\bibitem{power2017optimising}
P. Power, H. Rudd, J. Levitt ve F. Ensum,
``Optimising the press: Analysing the pressing strategies of the top five leagues
in European football,''
\textit{MIT Sloan Sports Analytics Conference}, 2017.

\bibitem{rathke2017expected}
A. Rathke,
``An examination of expected goals and shot quality from a goalie's perspective,''
\textit{Journal of Human Sport and Exercise}, vol. 12, no. 2, pp. 514--529, 2017.

\bibitem{statsbomb2023open}
StatsBomb,
``StatsBomb Open Data,''
[Çevrimiçi]. Erişim: \url{https://github.com/statsbomb/open-data}, 2023.

\bibitem{hamilton2020graph}
W. Hamilton,
\textit{Graph Representation Learning}.
Morgan \& Claypool Publishers, 2020.

\bibitem{kipf2017semi}
T. N. Kipf ve M. Welling,
``Semi-Supervised Classification with Graph Convolutional Networks,''
\textit{arXiv preprint arXiv:1609.02907}, 2017.

\bibitem{xu2019powerful}
K. Xu, W. Hu, J. Leskovec ve S. Jegelka,
``How Powerful are Graph Neural Networks?,''
\textit{arXiv preprint arXiv:1810.00826}, 2019.

\bibitem{decroos2019actions}
T. Decroos, L. Bransen, J. Van Haaren ve J. Davis,
``Actions Speak Louder than Goals: Valuing Player Actions in Football,''
\textit{ACM KDD}, 2019.

\end{thebibliography}

\end{document}
